%\documentclass[a4paper,11pt]{article}
\documentclass{article}
%%%%%%%%%%%%%%%%%%%%%%%%%%%%%%%%%%%%%%%%%%%%%%%%%%%%%%%%%%%%%%%%%%%%%%%%%%%%%%%%%%%%%%%%%%%%%%%%%%%%%%%%%%%%%%%%%%%%%%%%%%%%
\usepackage{graphics,graphicx}
\usepackage{amsmath,amssymb,graphics,graphicx}
\usepackage[ansinew]{inputenc}
\usepackage[usenames,dvipsnames]{color}


\usepackage{natbib}

\bibpunct{(}{)}{;}{a}{,}{,}

\textheight 24cm \textwidth 17cm \topmargin-2cm
%% \evensidemargin   -0.25cm
\oddsidemargin-0.2cm
%\pagestyle{empty}
\renewcommand{\baselinestretch}{1}

\begin{document}

\title{T�TULO}

\author{{Nombre Apellidos}\\
{\small Afiliaci�n..., Universidad Polit\'ecnica de Madrid, Spain}}

\date{}
\maketitle

%\title{}

%\address{}

\begin{abstract} Resumen breve. Resumen breve. Resumen breve. Resumen breve. Resumen breve. Resumen breve. Resumen breve. Resumen breve. Resumen breve. Resumen breve. Resumen breve.
\end{abstract}


\section{Introducci�n}

Al problema.



\section{Estado del arte}

Describir lo que han hecho otros, citando sus trabajos. Hay diferentes formatos dependiendo del tipo de trabajo: un art�culo como \cite{uncu2007}, un libro como \cite{hosmer-lemeshow2000}, un cap�tulo de libro como \cite{wold75} o unas actas de congreso (\cite{shakhnarovich2001icml}). 

    \subsection{M�s espec�fico}

Quiz� hagan falta subsecciones.

\section{Conclusiones y l�neas futuras}

Propias 



\bibliographystyle{plainnat}
\begin{thebibliography}{}

\bibitem[Hosmer and Lemeshow, 2000]{hosmer-lemeshow2000}
Hosmer, D.W. and Lemeshow, S. (2000) \textit{Applied Logistic Regression}. 2nd edn, J. Wiley and Sons, New York.

\bibitem[Shakhnarovich \textit{et al.}, 2001]{shakhnarovich2001icml}
Shakhnarovich, G., El-Yaniv, R. and Baram, Y. (2001) Smoothed bootstrap and statistical data cloning for classifier
evaluation. In \textit{Proceedings of the 18th International Conference on Machine Learning (ICML)}, Morgan Kaufmann, pp.
521--528.

\bibitem[Uncu and T\"urksen, 2007]{uncu2007}
Uncu, O. and T\"urksen, I.B. (2007) A novel feature selection
approach: Combining feature wrappers and filters. \textit{Information Sciences}, 177, 449--466.

\bibitem[Wold, 1975]{wold75} Wold, H. (1975) Soft modelling by latent variables: The non-linear iterative partial least squares ({NIPALS})
approach. In Gani, J. (ed), \textit{Perspectives in Probability and Statistics}, Academic Press, pp. 117--142.


\end{thebibliography}


\end{document}
